% Generated draft: requires mathematical and source review.
\documentclass[11pt]{article}
\usepackage[margin=25mm]{geometry}
\usepackage{fontspec}
\setmainfont{Latin Modern Roman}
\usepackage{amsmath,amssymb,mathtools,unicode-math}
\setmathfont{Latin Modern Math}
\usepackage{enumitem,needspace,xurl,hyperref}
\setlength{\parindent}{0pt}
\setlength{\parskip}{5pt}
\setlength{\emergencystretch}{3em}
\newcommand{\N}{\mathbb{N}}
\newcommand{\Z}{\mathbb{Z}}
\newcommand{\Q}{\mathbb{Q}}
\newcommand{\R}{\mathbb{R}}
\newcommand{\C}{\mathbb{C}}
\newcommand{\F}{\mathbb{F}}
\newcommand{\A}{\mathbb{A}}
\newcommand{\PP}{\mathbb{P}}
\newcommand{\E}{\mathbb{E}}
\newcommand{\eps}{\varepsilon}
\newcommand{\dd}{\,\mathrm{d}}
\begin{document}
\section*{UnsolvedMath ID: 30002912}
\textbf{Source problem code:} catalogue entry
\subsection{Definitions and mathematical statement}
Let \(X\) be a compact Kähler manifold of complex dimension \(n\), let \(\omega\) be its Kähler form, and let \(L\to X\) be a holomorphic line bundle. Its Chern character is \(\operatorname{ch}(L)=\exp(c_1(L))\). For an irreducible holomorphic subvariety \(V\subseteq X\) of complex dimension \(d\), define
\[
Z_V(L,\omega)=\int_{[V]}e^{-i[\omega]}\operatorname{ch}(L)
=\frac{1}{d!}\int_{[V]}\bigl(c_1(L)-i[\omega]\bigr)^d.
\]
Here integration means pairing the degree-\(2d\) component with the fundamental cycle of \(V\), including when \(V\) is singular.

For clarity about the named differential equation, one conventional normalization associates to a smooth Hermitian metric \(h\) on \(L\) the real curvature form \(\alpha_h=\frac{i}{2\pi}F_h\), representing \(c_1(L)\). With a constant phase \(\widehat\theta\in\R\), the deformed Hermitian--Yang--Mills equation in this normalization is
\[
\operatorname{Im}\!\left(e^{-i\widehat\theta}(\omega+i\alpha_h)^n\right)=0,
\qquad
\operatorname{Re}\!\left(e^{-i\widehat\theta}(\omega+i\alpha_h)^n\right)>0,
\]
where positivity means positivity as a real top-degree form. The supplied record does not specify its curvature normalization, prescribed phase or lift of the phase, or any additional admissibility condition; this displayed convention does not resolve those omissions.

The supplied question asks whether
\[
\text{the deformed Hermitian--Yang--Mills equation on }L\text{ admits a solution}
\quad\Longleftrightarrow\quad
\operatorname{Arg}Z_V(L,\omega)>-\operatorname{Arg}Z_X(L,\omega)
\quad\text{for every irreducible holomorphic subvariety }V\subset X.
\]
This is the recorded target, subject to unresolved qualifications: a real-valued branch or compatible lifts of the arguments must be specified, and each argument requires its central charge to be nonzero. The source also leaves unclear whether \(V\) must be proper, whether zero-dimensional subvarieties are included, and whether restrictions on dimension, phase, or positivity are intended. An ordered inequality between arguments is not defined solely by their values modulo \(2\pi\). Consequently, the supplied formulation does not determine a unique precise conjecture without these choices; no missing hypothesis or branch is imposed here.
\subsection{Short English statement}
The recorded question asks whether existence of a deformed Hermitian--Yang--Mills metric on a holomorphic line bundle over a compact Kähler manifold is equivalent to the stated argument inequality for every irreducible holomorphic subvariety. The phase convention, argument branches, nonvanishing requirements, and range of subvarieties need clarification before this becomes an unambiguous mathematical assertion.
\subsection{Sources}
\begin{itemize}
\item Local UnsolvedMath registry: supplied complete record for permanent ID \(30002912\), source problem code \texttt{catalogue entry}, titled ``Numerical Criterion for Deformed Hermitian--Yang--Mills Metrics.'' The collection attributes the record to Oberwolfach workshop open-problem discussions; no individual proposer is supplied.
\item Supplied report citation: \emph{Differentialgeometrie im Großen (2016)}, Oberwolfach Reports, pp.~1759--1807, DOI \texttt{10.4171/owr/2015/31}. Original report link: \url{https://doi.org/10.4171/owr/2015/31}. The record separately gives proposed year \(2015\); the apparent year discrepancy requires checking.
\item Supplied literature reference, labeled ``Collins Jacob and Yau primary paper'': \url{https://doi.org/10.4310/CJM.2020.V8.N2.A4}.
\item Supplied literature reference, labeled ``Lin primary paper'': \url{https://arxiv.org/abs/2201.01438}.
\item Supplied literature assessment dated exactly \texttt{2026-08-21}: status ``partially solved''; the record states that the existence criterion is verified on surfaces and under important additional hypotheses, while the full compact Kähler higher-dimensional conjecture remains unresolved. This is supplied dated curation, described in the record as literature triage derived from the supplied OWR curation, and has not been independently verified here. Time-sensitive status and cited claims require independent verification.
\end{itemize}
\end{document}
